% SEGMENT OLP-0703-S01
% Part: proof-theory
% Chapter: sequent-calculus
% Section: quantifiers

\documentclass[../../../include/open-logic-section]{subfiles}

\begin{document}

\olfileid{pt}{seq}{qua}

\olsection{ஒழுங்கான நிறுவல்களும் பதிலீடும்}

\LeftR{\lexists} மற்றும்~\RightR{\lforall}
விதிகளுக்குரிய ``தனிமாறி நிபந்தனைகள்''
காரணமாக அளவையடை விதிகளில் சில சிக்கல்கள்
ஏற்படுகின்றன. இவ்வனுமானங்களில் வரும்
!!{constant}~$a$ ஐ மீண்டும் பயன்படுத்தாமல்
இருந்தால், இவற்றில் பெரும்பாலானவற்றைச்
சமாளிக்கலாம். இதற்குப் பின்வரும்
வரையறையை அறிமுகப்படுத்துவோம்.

\begin{defn}
\Log{G1c} அல்லது \Log{G3c} இல்
!!^a{proof}~$\pi$ பின்வரும் இரு
நிபந்தனைகளையும் நிறைவேற்றினால் அது
\emph{ஒழுங்கானது}:
\begin{enumerate}
  \item எந்தத் தனிமாறி~$a$ உம் ஒன்றுக்கு மேற்பட்ட
  \LeftR{\lexists} அல்லது~\RightR{\lforall}
  அனுமானங்களில் தனிமாறியாகப் பயன்படுத்தப்படாது;
  \item $a$ ஒரு \LeftR{\lexists} அல்லது
  ~\RightR{\lforall} அனுமானத்தின் தனிமாறி எனில்,
  அது $\pi$ இல் அந்த அனுமானத்திற்கு மேலே
  மட்டுமே வரும்.
\end{enumerate}
\end{defn}

% SEGMENT OLP-0703-S02
\begin{lem}\ollabel{lem:subst-regular}
$\pi$ ஒழுங்கானதாகவும்,
$\Gamma \Sequent \Delta$ இல் முடிவதாகவும்,
$c$ என்பது $\pi$ இல் தனிமாறியாகப்
பயன்படாத !!a{constant} ஆகவும்,
$t$ என்பது $\pi$ இன் எந்தத் தனிமாறியையும்
கொண்டிராத சொல்லாகவும் இருந்தால்,
$\pi$ முழுவதிலும் $c$ வரும் ஒவ்வோர் இடத்திலும்
$t$ ஐப் பதிலிட்டுக் கிடைக்கும்
$\Subst{\pi}{t}{c}$ ஒரு ஒழுங்கான !!{proof} ஆகும்.
$\Subst{\pi}{t}{c}$ இன் இறுதித் தொடரணி
$\Subst{\Gamma}{t}{c}
\Sequent \Subst{\Delta}{t}{c}$;
இங்கு $\Subst{\Gamma}{t}{c} =
\Setabs{!A(t)}{!A(c) \in \Gamma}$.
\end{lem}

% SEGMENT OLP-0703-S03
\begin{proof}
  $\pheight{\pi}$ மீது தொகுத்தறிதல் செய்வோம்.
  $\pheight{\pi} = 0$ எனில் அது ஓர்
  அடிக்கோள் மட்டுமே. அப்போது
  $\Subst{\pi}{t}{c}$ உம் ஓர் அடிக்கோள்;
  ஏனெனில் அதிலுள்ள எந்த
  !!{formula}~$!A(c)$ உம் $!A(t)$ ஆக மாறும்.

  $\pheight{\pi} > 0$ எனில் $\pi$ இன்
  கடைசி அனுமானத்தைப் பொறுத்து நிலைகளைப்
  பிரிப்போம். கட்டமைப்பு விதிகள்
  (\Log{G1c} இல்) மற்றும் கூற்றிணைப்பிகளின்
  தருக்க விதிகள் தொடர்பான நிலைகள் எளியவை;
  அவை பயிற்சிகளாக விடப்படுகின்றன.
  $\pi$ ஓர் அளவையடை அனுமானத்தில் முடியும்
  நிலைகளைப் பார்ப்போம்.
  \begin{enumerate}

% SEGMENT OLP-0703-S04
    \item கடைசி அனுமானம் $\RightR{\lforall}$
    எனில், $\pi$ இவ்வடிவில் இருக்கும்:
    \[
    \AxiomC{}
    \RightLabel{$\pi'$}
    \Deduce$\Gamma(c) \fCenter \Delta'(c), !A(a,c)$
    \RightLabel{\RightR{\lforall}}
    \UnaryInf$\Gamma(c) \fCenter \Delta'(c), \lforall[x][!A(x,c)]$
    \DisplayProof
    \]
    தொகுத்தறிதல் கருதுகோளால்
    $\Subst{\pi'}{t}{c}$ என்பது
    $\Gamma(t) \fCenter \Delta'(t), !A(a,t)$
    என்பதற்கான ஒழுங்கான !!{proof}.
    $a$ என்பது $\pi$ இன் தனிமாறி என்பதால்,
    $t$ இல் $a$ வராது.
    எனவே $\Subst{\pi}{t}{c}$ இன்
    கடைசி அனுமானம்,
    \[
    \AxiomC{}
    \RightLabel{$\Subst{\pi'}{t}{c}$}
    \Deduce$\Gamma(t) \fCenter \Delta'(t), !A(a,t)$
    \RightLabel{\RightR{\lforall}}
    \UnaryInf$\Gamma(t) \fCenter \Delta'(t), \lforall[x][!A(x,t)]$
    \DisplayProof
    \]
    சரியானது: அதன் முடிவில் $a$ வராது.

% SEGMENT OLP-0703-S05
    \item கடைசி அனுமானம் \Log{G3c} இன்
    $\LeftR{\lforall}$ எனில்,
    $\pi$ இவ்வடிவில் இருக்கும்:
    \[
    \AxiomC{}
    \RightLabel{$\pi'$}
    \Deduce$!A(s(c),c), \lforall[x][!A(x,c)], \Gamma(c) \fCenter \Delta'(c)$
    \RightLabel{\LeftR{\lforall}}
    \UnaryInf$\lforall[x][!A(x,c)], \Gamma(c) \fCenter \Delta'(c)$
    \DisplayProof
    \]
    தொகுத்தறிதல் கருதுகோளால்,
    $\Subst{\pi'}{t}{c}$ என்பது
    $!A(s(t),t), \lforall[x][!A(x,t)], \Gamma(t) \fCenter
    \Delta'(t)$ என்பதற்கான ஒழுங்கான
    !!{proof}. $\Subst{\pi}{t}{c}$ இன்
    கடைசி அனுமானம்,
    \[
    \AxiomC{}
    \RightLabel{$\Subst{\pi'}{t}{c}$}
    \Deduce$!A(s(t),t), \lforall[x][!A(x,t)], \Gamma(t) \fCenter \Delta'(t)$
    \RightLabel{\LeftR{\lforall}}
    \UnaryInf$\lforall[x][!A(x,t)], \Gamma(t) \fCenter \Delta'(t)$
    \DisplayProof
    \]
    சரியானது. \Log{G1c} இன்
    \LeftR{\lforall} நிலையும் இதைப் போன்றதே;
    ஆனால் சற்றுக் குறைந்த சிக்கலுடையது.

% SEGMENT OLP-0703-S06
    \item கடைசி அனுமானம்
    \RightR{\lexists} அல்லது
    \LeftR{\lexists} எனில்: பயிற்சி.
  \end{enumerate}
  ஒவ்வொரு நிலையிலும் ஓர் ஒழுங்கான
  !!{proof} இன் முடிவில் ஒரு தொடரணியைச்
  சேர்த்தோம் (இரண்டு முன்கூற்றுகளுள்ள
  விதிகளில் இரண்டு ஒழுங்கான !!{proof}s இன்
  முடிவில்). புதிய இறுதித் தொடரணி,
  $\pi$ இன் இறுதித் தொடரணியிலிருந்து
  $c$ க்குப் பதிலாகச் சொல்~$t$
  வந்திருப்பதில் மட்டுமே வேறுபடுகிறது.
  $t$ இல் $\pi$ இன் தனிமாறிகள் வராததால்,
  $\pi$ மற்றும் $\Subst{\pi}{t}{c}$
  ஆகியவற்றின் தனிமாறிகள் ஒன்றே.
  ஒழுங்கான நிறுவலின் இறுதித் தொடரணியில்
  அதன் தனிமாறி வராததால்,
  புதிய இறுதித் தொடரணியிலும்
  $\Subst{\pi}{t}{c}$ இன் தனிமாறி வராது.
  ஆகவே $\Subst{\pi}{t}{c}$ ஒழுங்கானது.
\end{proof}

% SEGMENT OLP-0703-S07
\begin{prop}
$\pi$ என்பது \Log{G1c} அல்லது \Log{G3c}
இல் !!a{proof} எனில்,
அதே கட்டமைப்பும் (குறிப்பாக அதே !!{height} உம்)
உடைய ஒழுங்கான !!{proof}~$\pi'$ ஒன்று உள்ளது;
தனிமாறிகளை மாற்றியிருப்பதில் மட்டுமே
அது $\pi$ இலிருந்து வேறுபடுகிறது.
\end{prop}

% SEGMENT OLP-0703-S08
\begin{proof}
இந்த நிறுவலுக்குள் மட்டும், $\pi$ இன்
ஒரு தனிமாறி அனுமானத்தை
\emph{தூயது} என அழைப்போம், அதன் தனிமாறி
வேறொரு அனுமானத்திற்குத் தனிமாறியாகவும்
இருக்காமல், அந்த அனுமானத்தின் உடனடி
உட்-!!{proof} ஐத் தவிர $\pi$ இல்
வேறெங்கும் வராமலும் இருந்தால்;
இல்லையேல் அதை \emph{தூய்மையற்றது}
என்போம். $\pi$ இல் உள்ள தூய்மையற்ற
தனிமாறி அனுமானங்களின் எண்ணிக்கை $n$ என்க.
$n$ மீது தொகுத்தறிந்து $\pi$ ஐ
வேண்டிய ஒழுங்கான !!{proof}~$\pi'$ ஆக
மாற்றலாம் என்பதைக் காட்டுவோம்.
$n=0$ எனில் $\pi$ இன் எல்லாத் தனிமாறி
அனுமானங்களும் தூயவை; எனவே
$\pi$ ஒழுங்கானது, நிறுவ வேண்டியது ஏதுமில்லை.
இல்லையேல் மேலே வேறு தூய்மையற்ற
தனிமாறி அனுமானம் இல்லாத ஒன்றைத்
தேர்க; அது \RightR{\lforall} என்க.
அதில் முடியும் $\pi$ இன் உட்-!!{proof}~$\pi_1$
இவ்வடிவில் இருக்கும்:
    \[
    \AxiomC{}
    \RightLabel{$\pi_2$}
    \Deduce$\Gamma \fCenter \Delta, !A(c)$
    \RightLabel{\RightR{\lforall}}
    \UnaryInf$\Gamma \fCenter \Delta, \lforall[x][!A(x)]$
    \DisplayProof
    \]

% SEGMENT OLP-0703-S09
$c$ தனிமாறி என்பதால் $\Gamma$,
$\Delta$, $\lforall[x][!A(x)]$
ஆகியவற்றில் வராது. $\pi_2$ இல்
உள்ள தனிமாறி அனுமானங்கள் அனைத்தும்
தூயவை என்பதால் அது ஒழுங்கானது;
அதனுள் அந்தக் குறி தனிமாறியாகப் பயன்படுத்தப்படாது.
$\pi$ இல் வராத !!a{constant}~$a$
ஐ எடுக்கவும்.
\olref{lem:subst-regular} படி,
$\Subst{\pi_2}{a}{c}$ என்பது
$\Gamma \fCenter \Delta, !A(a)$
என்பதற்கான ஒழுங்கான நிறுவல்;
அதன் மீது $\RightR{\lforall}$
விதியைப் பயன்படுத்தலாம்.
$\pi$ இல் $\pi_1$ க்குப் பதிலாக
பின்வரும் நிறுவல் $\pi_1'$ ஐ வைத்தால்,
    \[
    \AxiomC{}
    \RightLabel{$\Subst{\pi_2}{a}{c}$}
    \Deduce$\Gamma \fCenter \Delta, !A(a)$
    \RightLabel{\RightR{\lforall}}
    \UnaryInf$\Gamma \fCenter \Delta, \lforall[x][!A(x)]$
    \DisplayProof
    \]
ஒரு தூய்மையற்ற தனிமாறி அனுமானம்
தூயதாக மாறுகிறது; வேறுபாடு தனிமாறி
$c$ க்குப் பதிலாக~$a$ வந்தது மட்டுமே.
கிடைத்த நிறுவலில் தூய்மையற்ற
தனிமாறி அனுமானங்கள் $n-1$ உள்ளன;
எனவே தொகுத்தறிதல் கருதுகோளால்
முடிவு கிடைக்கிறது.
\end{proof}

% SEGMENT OLP-0703-S10
இப்போது நிறுவிய முன்மொழிவைத்
தனியே மேற்கோளிட மாட்டோம். ஆனால்
நாம் கருதும் எல்லா !!{proof}s உம்
ஒழுங்கானவை எனக் கொள்வோம்: அவ்வாறு
இல்லையெனில், தனிமாறிகளை மாற்றி
அவற்றை ஒழுங்காக்கலாம்.

\end{document}
